\author{Lukas Denk and Melvatha R. Chee}
\title{Nine Navajo narratives}
%\subtitle{}
\renewcommand{\lsISBNdigital}{978-3-96110-580-9}
\renewcommand{\lsISBNhardcover}{978-3-98554-197-3}
\BookDOI{10.5281/zenodo.20522183}
\typesetter{Christian Döhler}
\proofreader{Matthew Korte}
\lsCoverTitleSizes{45pt}{15mm}
\renewcommand{\lsSeries}{otc}
\renewcommand{\lsSeriesNumber}{10}
\renewcommand{\lsID}{593}
\dedication{In memory of the narrators and for their descendants}

\renewcommand{\lsDedicationFont}{\fontsize{15pt}{8mm}\selectfont}
\BackBody{\textit{Éí bik'ehgo kééhwiit'į́}. `In accord with that, we live.' So speaks a medicine man after recounting how the four sacred mountains were set in place to anchor the world of the Diné. To live in accord with something is to know where you stand, even when the ground shifts. And the ground shifted constantly across the late 19th and early 20th centuries: at Round Rock where armed men faced down U.S. cavalry; in the canyons where boys fled from bears; in the corrals from which sheep were driven away by government agents; in the frozen winter when pilots dropped hay through fog to save starving communities. \textit{Nine Navajo Narratives} captures these moments through nine annotated texts: three firsthand accounts of the 1892 ``Trouble at Round Rock'' over forced boarding schools; a sacred creation story; a traditional tale explaining a mother-in-law custom; a testimony of relief flights during the blizzard of 1949; a vivid childhood narrative of adventure and taboo; and two fierce political critiques of the devastating livestock reduction policies of the 1930s.\\ Drawn from the foundational collections of Young and Morgan, all texts are presented in parallel translation with full interlinear glossing and an introductory grammar sketch—an essential resource for understanding Navajo history, language, and survival.}


